\documentclass[10pt, letterpaper]{article}

% Packages:
\usepackage[
    ignoreheadfoot, % set margins without considering header and footer
    top=0.5 cm, % seperation between body and page edge from the top
    bottom=1 cm, % seperation between body and page edge from the bottom
    left=0.8 cm, % seperation between body and page edge from the left
    right=0.8 cm, % seperation between body and page edge from the right
    footskip=1.0 cm, % seperation between body and footer
    % showframe % for debugging 
]{geometry} % for adjusting page geometry
\usepackage{titlesec} % for customizing section titles
\usepackage{tabularx} % for making tables with fixed width columns
\usepackage{array} % tabularx requires this
\usepackage[dvipsnames]{xcolor} % for coloring text
\definecolor{primaryColor}{RGB}{0, 0, 0} % define primary color
\usepackage{enumitem} % for customizing lists
\usepackage{fontawesome5} % for using icons
\usepackage{amsmath} % for math
\usepackage[
    pdftitle={Ahmed MAKROUM's CV},
    pdfauthor={Ahmed MAKROUM},
    pdfcreator={LaTeX with RenderCV},
    colorlinks=true,
    urlcolor=primaryColor
]{hyperref} % for links, metadata and bookmarks
\usepackage[pscoord]{eso-pic} % for floating text on the page
\usepackage{calc} % for calculating lengths
\usepackage{bookmark} % for bookmarks
\usepackage{lastpage} % for getting the total number of pages
\usepackage{changepage} % for one column entries (adjustwidth environment)
\usepackage{paracol} % for two and three column entries
\usepackage{ifthen} % for conditional statements
\usepackage{needspace} % for avoiding page brake right after the section title
\usepackage{iftex} % check if engine is pdflatex, xetex or luatex

% Ensure that generate pdf is machine readable/ATS parsable:
\ifPDFTeX
    \input{glyphtounicode}
    \pdfgentounicode=1
    \usepackage[T1]{fontenc}
    \usepackage[utf8]{inputenc}
    \usepackage{lmodern}
\fi

\usepackage{charter}

% Some settings:
\raggedright
\AtBeginEnvironment{adjustwidth}{\partopsep0pt} % remove space before adjustwidth environment
\pagestyle{empty} % no header or footer
\setcounter{secnumdepth}{0} % no section numbering
\setlength{\parindent}{0pt} % no indentation
\setlength{\topskip}{0pt} % no top skip
\setlength{\columnsep}{0.15cm} % set column seperation
\pagenumbering{gobble} % no page numbering

\titleformat{\section}{\needspace{4\baselineskip}\bfseries\large}{}{0pt}{}[\vspace{1pt}\titlerule]

\titlespacing{\section}{
    % left space:
    -1pt
}{
    % top space:
    0.2 cm
}{
    % bottom space:
    0.2 cm
} % section title spacing

\renewcommand\labelitemi{$\vcenter{\hbox{\small$\bullet$}}$} % custom bullet points
\newenvironment{highlights}{
    \begin{itemize}[
        topsep=0.10 cm,
        parsep=0.10 cm,
        partopsep=0pt,
        itemsep=0pt,
        leftmargin=0 cm + 10pt
    ]
}{
    \end{itemize}
} % new environment for highlights


\newenvironment{highlightsforbulletentries}{
    \begin{itemize}[
        topsep=0.10 cm,
        parsep=0.10 cm,
        partopsep=0pt,
        itemsep=0pt,
        leftmargin=10pt
    ]
}{
    \end{itemize}
} % new environment for highlights for bullet entries

\newenvironment{onecolentry}{
    \begin{adjustwidth}{
        0 cm + 0.00001 cm
    }{
        0 cm + 0.00001 cm
    }
}{
    \end{adjustwidth}
} % new environment for one column entries

\newenvironment{twocolentry}[2][]{
    \onecolentry
    \def\secondColumn{#2}
    \setcolumnwidth{\fill, 4.5 cm}
    \begin{paracol}{2}
}{
    \switchcolumn \raggedleft \secondColumn
    \end{paracol}
    \endonecolentry
} % new environment for two column entries

\newenvironment{threecolentry}[3][]{
    \onecolentry
    \def\thirdColumn{#3}
    \setcolumnwidth{, \fill, 4.5 cm}
    \begin{paracol}{3}
    {\raggedright #2} \switchcolumn
}{
    \switchcolumn \raggedleft \thirdColumn
    \end{paracol}
    \endonecolentry
} % new environment for three column entries

\newenvironment{header}{
    \setlength{\topsep}{0pt}\par\kern\topsep\centering\linespread{1.5}
}{
    \par\kern\topsep
} % new environment for the header

\newcommand{\placelastupdatedtext}{% \placetextbox{<horizontal pos>}{<vertical pos>}{<stuff>}
  \AddToShipoutPictureFG*{% Add <stuff> to current page foreground
    \put(
        \LenToUnit{\paperwidth-2 cm-0 cm+0.05cm},
        \LenToUnit{\paperheight-1.0 cm}
    ){\vtop{{\null}\makebox[0pt][c]{
        \small\color{gray}\textit{Last updated in September 2024}\hspace{\widthof{Last updated in September 2024}}
    }}}%
  }%
}%

% save the original href command in a new command:
\let\hrefWithoutArrow\href

% new command for external links:


\begin{document}
    \newcommand{\AND}{\unskip
        \cleaders\copy\ANDbox\hskip\wd\ANDbox
        \ignorespaces
    }
    \newsavebox\ANDbox
    \sbox\ANDbox{$|$}

    \begin{header}
        \fontsize{25 pt}{25 pt}\selectfont Ahmed MAKROUM

        \vspace{5 pt}

        \normalsize
        \mbox{Casablanca, Maroc}%
        \kern 5.0 pt%
        \AND%
        \kern 5.0 pt%
        \mbox{\hrefWithoutArrow{mailto:ahmedmakroum3@gmail.com}{ahmedmakroum3@gmail.com}}%
        \kern 5.0 pt%
        \AND%
        \kern 5.0 pt%
        \mbox{\hrefWithoutArrow{tel:+212 664716219}{06 64 71 62 19}}%
        \kern 5.0 pt%
        \AND%
        \kern 5.0 pt%
        \mbox{\hrefWithoutArrow{https://makroum.website/}{makroum.website}}%
        \kern 5.0 pt%
        \AND%
        \kern 5.0 pt%
        \mbox{\hrefWithoutArrow{https://www.linkedin.com/in/ahmed-makroum/}{/in/ahmed-makroum}}%
        \kern 5.0 pt%
        \AND%
        \kern 5.0 pt%
        \mbox{\hrefWithoutArrow{https://github.com/ahmedmakroum}{github.com/ahmedmakroum}}%
    \end{header}

    \vspace{5 pt - 0.3 cm}


    \section{Profil}



        
        \begin{onecolentry}
            Ingénieur logiciel développant des applications full-stack, des systèmes backend et des pipelines de données.

            Expérience sur .NET, React et des outils modernes d’ingénierie data, avec livraison de solutions de bout en bout, de l’architecture à la production. Axé sur la performance, la conception de systèmes clairs et la construction d’applications fiables utilisées en conditions réelles.
        \end{onecolentry}





    
    \section{Expérience}


     \begin{twocolentry}{
            Fév. 2025 –  Présent
        }
            \textbf{Ingénieur logiciel}, Attijariwafa bank \end{twocolentry}

        \vspace{0.10 cm}
        \begin{onecolentry}
            \begin{highlights}
\item Développement d’applications full-stack avec .NET et React pour soutenir les workflows métier et améliorer l’efficacité opérationnelle.
    \item Intégration des API Microsoft Graph pour centraliser les données utilisateurs, permissions et collaboration entre systèmes.
    \item Développement de web parts SPFx (React) dans SharePoint, améliorant l’accessibilité et l’adoption des outils internes.
    \item Conception et mise en place d’architectures backend (API REST, authentification, contrôle d’accès par rôles) garantissant évolutivité et maintenabilité.



            \end{highlights}
        \end{onecolentry}


        \vspace{0.2 cm}


    \begin{twocolentry}{
            Sept. 2025 –  Déc. 2025
        }
            \textbf{Ingénieur Data}, Metaverse \end{twocolentry}

        \vspace{0.10 cm}
        \begin{onecolentry}
            \begin{highlights}
                \item Conception d’architectures data modulaires sur AWS, GCP et environnements on-premise, standardisant les déploiements et réduisant la complexité de mise en place.
                \item Automatisation du provisionnement d’infrastructure avec Terraform, permettant des environnements évolutifs et reproductibles en quelques minutes.


            \end{highlights}
        \end{onecolentry}


        \vspace{0.2 cm}
        
        \begin{twocolentry}{
            Mars 2025 – Sept. 2025
        }
            \textbf{Stage – Ingénieur Data}, Allianz Maroc \end{twocolentry}

        \vspace{0.10 cm}
        \begin{onecolentry}
            \begin{highlights}
\item Construction de pipelines de données (NiFi + Spark) consolidant des données multi-sources dans un entrepôt PostgreSQL centralisé.
    \item Automatisation des extractions de données, réduisant le temps de livraison des rapports de plusieurs heures à quelques minutes.
    \item Mise en place de tableaux de bord self-service (Metabase), permettant aux équipes métier de suivre les KPI sans dépendance technique.
    \item Développement d’une application basée sur des microservices (Spring Boot + Next.js) pour la gestion des sinistres, utilisée par 50+ utilisateurs quotidiens.

            \end{highlights}
        \end{onecolentry}


        \vspace{0.2 cm}
        
        \begin{twocolentry}{
            Juin 2024 – Août 2024
        }
            \textbf{Stage – Ingénieur Data}, Boti School \end{twocolentry}

        \vspace{0.10 cm}
        \begin{onecolentry}
            \begin{highlights}
\item Développement de pipelines Apache Beam (GCP) traitant des données d’usage à grande échelle pour un suivi précis de la consommation.
    \item Conception d’un entrepôt BigQuery supportant des systèmes automatisés de facturation par utilisateur.
    \item Mise en place de tableaux de bord en temps réel (Looker Studio), supprimant le reporting financier manuel.
            \end{highlights}
        \end{onecolentry}



        \vspace{0.2 cm}

        \begin{twocolentry}{
            Juin 2024 – Août 2024
        }
            \textbf{Stage – Développeur Web et Mobile}, 6Solutions\end{twocolentry}

        \vspace{0.10 cm}
        \begin{onecolentry}
            \begin{highlights}
\item Développement d’une plateforme multi-services centralisant des services juridiques, médicaux et financiers.
    \item Développement d’une architecture frontend/backend évolutive avec Angular et Spring Boot.
    \item Livraison d’une application mobile multiplateforme avec Flutter.
    
            \end{highlights}
        \end{onecolentry}



        \vspace{0.2 cm}

        \begin{twocolentry}{
            Juil. 2023 – Août 2023
        }
            \textbf{Stage – Développeur Jeu Vidéo}, Finso\end{twocolentry}

        \vspace{0.10 cm}
        \begin{onecolentry}
            \begin{highlights}
                \item Développement d’un jeu éducatif introduisant des notions de gestion et de stratégie à travers un gameplay interactif.
                \item Réalisation du projet sous Unity avec des mécaniques interactives et une expérience utilisateur optimisée pour un jeune public.
            \end{highlights}
        \end{onecolentry}
    \section{Formation}



        
        \begin{twocolentry}{
            Sept 2020 – Mars 2025
        }
            \textbf{EMSI – École Marocaine des Sciences de l’Ingénieur}, Diplôme d’Ingénieur en Informatique et Réseaux option MIAGE \end{twocolentry}




    
    \section{Compétences}



        
        \begin{onecolentry}
            \textbf{Stack principal:} .NET, React, Next.js, Spring Boot

            \vspace{0.1 cm}

            \textbf{Ingénierie Data:} Spark, NiFi, Airflow, dbt, Kafka

            \vspace{0.1 cm}

            \textbf{Cloud \& Infrastructure:} AWS, GCP, Docker, Terraform


        \end{onecolentry}


    

\end{document}
